% GENERATED_BY: oc_core_monograph_translation_draft_pipeline_v1.py
% AI_ASSISTED_TRANSLATION_DRAFT: TRUE
% THEOREM_REVIEW_STATUS: PENDING
% THEOREM_REVIEW_MODE: FORMAL_AI_GATE__CODEX_CLI_CHATGPT
% THEOREM_REVIEW_REVIEWER_ID:
% THEOREM_REVIEW_AUTHORITY_CLASS:
% THEOREM_REVIEWED_AT:
% THEOREM_REVIEW_DOSSIER_REF:
% SOURCE_LANGUAGE: EN
% TARGET_LANGUAGE: RU
% SOURCE_EN_REF: content/m_spaces/m11.tex
% ================================================================
% ==== FILE: content/m_spaces/m11.tex
% ================================================================

% ==============================
%  Ontology of Continua — Core
%  Meta-Space: M11
%  FULL MODULE — FINAL VERSION
% ==============================

\subsubsection{$M_{11}$ Обзор}
\label{sec:m11-overview}

$M_{11}$ — высшее мета-пространство, необходимое для континуумов до уровня $K_{11}$.
Если $M_{10}$ содержит категории категорий моделей (мета-теоретические структуры),
$M_{11}$ содержит \emph{иерархии мета-теоретических башен} и обеспечивает
глобальное логическое пространство для их совместимости, рефлексии
и стабилизации более высокого порядка.

В этом пространстве фундаментальные единицы — \emph{мета-мета-категориальные} объекты:
структурированные башни
\[
(\mathcal{C}^{(0)}, \mathcal{C}^{(1)}, \ldots, \mathcal{C}^{(n)}, \ldots)
\]
где каждый $\mathcal{C}^{(i)}$ сам является категорией категорий на уровне $i$,
а морфизмы между башнями — это преобразования, сохраняющие их полную
вертикальную структуру.

Следовательно, $M_{11}$ — минимальное логическое пространство, в котором
\emph{мета-иерархические континуумы} могут существовать.

% ---------------------------------------------------------------
\subsubsection*{1. Допустимое пространство конфигураций \texorpdfstring{$\Omega(M_{11})$}{\Omega(M_{11})}}

\[
\Omega(M_{11}) =
\left\{
\left( \mathbf{T},\ \mathfrak{F},\ A_{11},\
\Theta_{11},\ C_{11},\Sigma_{11} \right)
\ \middle|\
\text{выполняются условия глобальной иерархической когерентности}
\right\}.
\]

Где:

\begin{itemize}
    \item $\mathbf{T}$ — допустимые конечные или бесконечные башни мета-категорий,
    \item $\mathfrak{F}$ — мета-функторы, действующие на башнях,
    \item $A_{11}$ — оси вариативности мета-иерархий,
    \item $\Theta_{11}$ — глобальные пороги когерентности на верхнем уровне,
    \item $C_{11}$ — циклы транс-иерархической самоорганизации,
    \item $\Sigma_{11}$ — структурные ограничения, определяющие $M_{11}$.
\end{itemize}

Допустимость требует:

\begin{enumerate}
    \item иерархия должна замыкаться относительно вертикальной композиции,
    \item отсутствие расходимости когерентности между уровнями,
    \item существование хотя бы одного устойчивого транс-иерархического цикла,
    \item совместимость с $M_{10}$ через каноническую проекцию
          $\pi_{10}:M_{11}\to M_{10}$,
    \item конечность глобального мета-тензора натяжения.
\end{enumerate}

\[
\Omega(K_{11}) \subseteq \Omega(M_{11})
\quad\text{(Теорема 8: необходимость и достаточность совместимости)}.
\]

% ---------------------------------------------------------------
\subsubsection*{2. Оси $A(M_{11})$}

\[
A(M_{11}) =
\{
A_{\mathrm{tower}},\
A_{\mathrm{coherence}},\
A_{\mathrm{depth}},\
A_{\mathrm{reflection}},\
A_{\mathrm{translation}},\
A_{\mathrm{hierarchical\ consistency}},\
A_{\mathrm{global}}
\}.
\]

Интерпретация:

\begin{itemize}
    \item $A_{\mathrm{tower}}$ — вариативность по уровням мета-теоретической башни,
    \item $A_{\mathrm{coherence}}$ — поддержание когерентности на множестве уровней,
    \item $A_{\mathrm{depth}}$ — иерархическая глубина и структурная рекурсия,
    \item $A_{\mathrm{reflection}}$ — самописание по уровням,
    \item $A_{\mathrm{translation}}$ — перевод между различными мета-иерархиями,
    \item $A_{\mathrm{hierarchical\ consistency}}$ — сохранение структурных правил по всей башне,
    \item $A_{\mathrm{global}}$ — оси, управляющие глобальной допустимостью целых иерархий.
\end{itemize}

Ключевое свойство:

\[
\dim(M_{11}) > \dim(M_{10})
\quad\text{(Теорема: монотонность и необратимость роста размерности)}.
\]

% ---------------------------------------------------------------
\subsubsection*{3. Потенциалы \texorpdfstring{$P(M_{11})$}{P(M_{11})}}

\[
\begin{aligned}
P(M_{11}) = \{&
F_{\mathrm{global\ coherence}},\
F_{\mathrm{hierarchical\ depth}},\
F_{\mathrm{meta\mbox{-}reflection}},\\
&
F_{\mathrm{compatibility}},\
F_{\mathrm{translation}},\
F_{\mathrm{unification}},\
F_{\mathrm{stability}}
\}.
\end{aligned}
\]

Интерпретация:

\begin{itemize}
    \item $F_{\mathrm{global\ coherence}}$ — способность сохранять когерентность на всех уровнях,
    \item $F_{\mathrm{hierarchical\ depth}}$ — способность поддерживать глубокие иерархические структуры,
    \item $F_{\mathrm{meta\mbox{-}reflection}}$ — способность отражать целые иерархии,
    \item $F_{\mathrm{compatibility}}$ — согласованность между уровнями,
    \item $F_{\mathrm{translation}}$ — отображение между различными башнями,
    \item $F_{\mathrm{unification}}$ — унификация на разных иерархических масштабах,
    \item $F_{\mathrm{stability}}$ — общая стабильность динамики мета-иерархий.
\end{itemize}

Существование $K_{11}$ требует:

\[
F_{\mathrm{global\ coherence}} > \Theta_{11}^{\mathrm{collapse}},
\qquad
F_{\mathrm{meta\mbox{-}reflection}} > \Theta_{11}^{\mathrm{self\mbox{-}break}}.
\]

% ---------------------------------------------------------------
\subsubsection*{4. Пороги $\Theta(M_{11})$}

Критические пороги:

\begin{itemize}
    \item $\Theta_{11}^{\mathrm{coherence}}$ — минимальная глобальная когерентность,
    \item $\Theta_{11}^{\mathrm{depth}}$ — минимальная структурная глубина,
    \item $\Theta_{11}^{\mathrm{translation}}$ — порог интероперабельности,
    \item $\Theta_{11}^{\mathrm{collapse}}$ — граница, на которой иерархия даёт сбой,
    \item $\Theta_{11}^{\mathrm{self\mbox{-}break}}$ — нестабильность самоссылки.
\end{itemize}

Превышение $\Theta_{11}^{\mathrm{collapse}}$ вызывает иерархическую фрагментацию:

\[
\mathbf{T} \to (\mathcal{C}^{(i)}) \quad \text{без допустимой вертикальной структуры}.
\]

% ---------------------------------------------------------------
\subsubsection*{5. Потоки $J(M_{11})$}

\[
J(M_{11}) =
(
J_{\mathrm{tower}},\
J_{\mathrm{coherence}},\
J_{\mathrm{translation}},\
J_{\mathrm{meta\mbox{-}inference}},\
\nabla_i T_{11},\
J_{\mathrm{reflection}},\
J_{\mathrm{global}}
).
\]

Интерпретация:

\begin{itemize}
    \item $J_{\mathrm{tower}}$ — потоки, действующие на иерархическую глубину,
    \item $J_{\mathrm{coherence}}$ — потоки, восстанавливающие многоуровневую согласованность,
    \item $J_{\mathrm{translation}}$ — потоки, отображающие одну башню в другую,
    \item $J_{\mathrm{meta\mbox{-}inference}}$ — вывод на всём протяжении башни,
    \item $\nabla_i T_{11}$ — градиенты глобального мета-тензора натяжения,
    \item $J_{\mathrm{reflection}}$ — рекурсия и рефлексия по уровням,
    \item $J_{\mathrm{global}}$ — потоки, определяющие глобальную допустимость.
\end{itemize}

% ---------------------------------------------------------------
\subsubsection*{6. Циклы $C(M_{11})$}

Фундаментальные циклы:

\begin{itemize}
    \item $C_{\mathrm{hierarchy}}$ — формирование, уточнение и стабилизация башни,
    \item $C_{\mathrm{coherence}}$ — петли восстановления когерентности по уровням,
    \item $C_{\mathrm{translation}}$ — сопоставление башен и проверка эквивалентности,
    \item $C_{\mathrm{reflection}}$ — рекурсивные циклы, позволяющие
          понимать иерархию изнутри,
    \item $C_{\mathrm{global}}$ — глобальные циклы допустимости, обеспечивающие
          $\Omega(K_{11})\subseteq\Omega(M_{11})$.
\end{itemize}

Критерий жизни:

\[
\oint_{C_{\mathrm{global}}} dA_{\mathrm{global}} \approx 0.
\]

% ---------------------------------------------------------------
\subsubsection*{7. Граница \texorpdfstring{$\partial\Omega(M_{11})$}{\partial\Omega(M_{11})}}

Приближение к границе:

\begin{itemize}
    \item иерархическая дивергенция,
    \item потеря когерентности между уровнями,
    \item нерешаемые противоречия,
    \item разгон рекурсии,
    \item нарушение глобальной допустимости.
\end{itemize}

Пересечение $\partial\Omega(M_{11})$ принуждает к свёртыванию в структуры, допустимые в $M_{10}$.

% ---------------------------------------------------------------
\subsubsection*{8. Связь с $M_{10}$ и высшими мета-пространствами}

\textbf{Связь с $M_{10}$:}

\[
A_{\mathrm{hierarchical\ consistency}} \in A(M_{11}) \setminus A(M_{10}),
\]
следовательно $M_{11}$ строго расширяет $M_{10}$.

Вертикальная проекция:

\[
\pi_{10}: M_{11} \to M_{10}
\]
сбрасывает иерархическую глубину и оставляет только мета-теоретическую структуру.

\textbf{К $M_{12}$ (если определено):}

$M_{12}$ должен содержать универсальные мета-мета-логические структуры и
глобальные ограничения на целые семейства мета-иерархий.

$M_{11}$ — минимальное пространство, в котором такая расширяемость становится допустимой.
